\documentclass[11pt]{article}
\usepackage{geometry}
\geometry{a4paper, margin=2.5cm}
\usepackage{amsmath, amssymb, amsthm, array, booktabs}
\usepackage{parskip}
\usepackage{hyperref}
\hypersetup{colorlinks=true, linkcolor=blue, citecolor=blue, urlcolor=blue}

\newtheorem{theorem}{Theorem}
\newtheorem{proposition}{Proposition}
\newtheorem{remark}{Remark}

\title{Disproof of conjecture \texttt{00000001068}}
\author{earthking11}
\date{2026-10-01}

\begin{document}
\maketitle

\section{The conjecture}

Conjecture \texttt{00000001068} of The Last Math Competition states:

\begin{quote}
\textbf{English.} Definition: A sum-free set ($A + A$ disjoint from $A$).
Conjecture: The maximal sum-free subsets of $\mathbb{F}_p$ are exactly the
intervals $((p+1)/3,\, 2(p-1)/3)$, uniquely up to dilation (the complete
Diananda--Yap characterization over finite fields).
\end{quote}
(The original also gives a Chinese version; it is quoted in \texttt{README.md}.)

The parenthetical reference is to the classical Diananda--Yap theorem, which
indeed characterizes the \emph{maximum-cardinality} sum-free subsets of
$\{1,\dots,n\}$ (and, with the right endpoints, of $\mathbb{F}_p$). The
conjecture as written, however, gets the interval endpoints wrong, and is
false already at $p = 11$.

\section{The counterexample}

Work with integer representatives $1, \dots, p-1$ of $\mathbb{F}_p^\times$
(the interval never contains $0$ for the primes below, and adding $0$ to a
sum-free set always destroys sum-freeness, so this loses no generality).

For $p = 11$ the conjectured interval is, with strict inequalities,
\[
\left(\frac{11+1}{3},\, \frac{2(11-1)}{3}\right) = \left(4,\, \frac{20}{3}\right),
\]
i.e.\ the integer set
\[
I_{11} = \{5, 6\}, \qquad |I_{11}| = 2.
\]

\begin{theorem}
$S = \{1, 3, 8, 10\} \subseteq \mathbb{F}_{11}$ is sum-free, has $4$ elements,
and is maximal by inclusion. In particular $S$ is not a dilation of $I_{11}$,
so the conjecture is false at $p = 11$.
\end{theorem}

\begin{proof}
The $16$ ordered pair-sums $a + b \pmod{11}$ with $a, b \in S$ are
\[
\begin{array}{c|cccc}
+      & 1 & 3 & 8 & 10 \\ \hline
1      & 2 & 4 & 9 & 0  \\
3      & 4 & 6 & 0 & 2  \\
8      & 9 & 0 & 5 & 7  \\
10     & 0 & 2 & 7 & 9
\end{array}
\]
so $S + S = \{0, 2, 4, 5, 6, 7, 9\}$, disjoint from $S = \{1, 3, 8, 10\}$:
$S$ is sum-free. Since dilation $x \mapsto c x$ with $c \ne 0$ is a bijection
of $\mathbb{F}_{11}$, every dilation of the $2$-element interval $I_{11}$ still
has $2$ elements, whereas $|S| = 4$; hence $S$ cannot be a dilation of $I_{11}$.

Moreover $S$ is maximal by inclusion: for each of the $7$ points
$x \in \mathbb{F}_{11} \setminus S = \{0, 2, 4, 5, 6, 7, 9\}$, the set
$S \cup \{x\}$ fails to be sum-free (e.g.\ $0 + 1 = 1$, $2 + 8 = 10$,
$4 + 7$ is not needed---each case is a finite check; see
\texttt{reproduce.py}). Thus $S$ is a ``maximal sum-free set'' in the
inclusion sense that is not an interval dilation, refuting the conjecture
under either reading of ``maximal''.
\end{proof}

\begin{remark}[The case $p = 5$]
At $p = 5$ the conjecture's interval $\bigl(\frac{5+1}{3}, \frac{2(5-1)}{3}\bigr)
= (2, 8/3)$ is \emph{empty}, while $\{1, 4\} \subseteq \mathbb{F}_5$ is a
nonempty sum-free set ($1+1=2$, $1+4=0$, $4+4=3 \pmod 5$). So the ``interval''
cannot be the maximum sum-free set either.
\end{remark}

\section{Formalization}

The finite checks above are formalized in Lean~4 (see \texttt{lean4/}):
\texttt{Main.lean} uses core Lean only (\texttt{import Std}), and every theorem
is proved by \texttt{decide} (kernel computation over \texttt{Nat}; the
rational bound $x < 20/3$ is encoded as $3x < 20$). \texttt{Check.lean} runs
\texttt{\#print axioms} on all twelve theorems: none depends on any axiom
(no \texttt{sorryAx}, no \texttt{Lean.ofReduceBool}, no Mathlib).

\section{Conclusion}

Conjecture \texttt{00000001068} is false: at $p = 11$ the set
$\{1, 3, 8, 10\}$ is a $4$-element maximal-by-inclusion sum-free subset of
$\mathbb{F}_{11}$, strictly larger than the conjectured $2$-element interval
$\{5, 6\}$ and not a dilation of it. The AI-generated conjecture misstates
the interval endpoints of the Diananda--Yap characterization.

\end{document}
